%-----------------------------------------------------------------------
\subsection{Gauge Freedom (Physics)}
\label{sec:instance-gauge-theory}
%-----------------------------------------------------------------------

Gauge invariance is a foundational principle of modern physics: physically
measurable quantities are invariant under gauge transformations, but the
mathematical description of the field is not
\citep{jackson1999}.  We formalize a discrete $\mathbb{Z}_2$ gauge
theory on a triangle graph (three vertices, three edges) to exhibit the
Explanation Impossibility in physics.

\paragraph{The explanation system.}
Define $\mathcal{S}_{\mathrm{gauge}} = (\Params, \sH, Y, \mathsf{obs},
\mathsf{exp}, \incomp)$ as follows:
\begin{itemize}
  \item $\Params = \sH = \texttt{EdgeConfig}$: each configuration assigns a
    $\mathbb{Z}_2$ (Boolean) label to each of the three edges
    $e_{01}, e_{12}, e_{02}$;
  \item $Y = \texttt{Bool}$: the \emph{holonomy} around the triangle,
    $h(g) = e_{01} \oplus e_{12} \oplus e_{02}$, the discrete analogue of
    the Wilson loop;
  \item $\mathsf{obs} = \texttt{holonomy}$;
  \item $\mathsf{exp} = \mathrm{id}$ (the full edge configuration is the
    ``explanation'');
  \item $\incomp(g_1, g_2) \iff g_1 \ne g_2$.
\end{itemize}

A \emph{gauge transform at vertex~0} flips the two incident edges
($e_{01}$ and $e_{02}$), leaving $e_{12}$ unchanged.

\paragraph{The Rashomon property.}
The witnesses are $g_1 = (\mathtt{true}, \mathtt{false}, \mathtt{false})$
and $g_2 = (\mathtt{false}, \mathtt{false}, \mathtt{true})$.  They satisfy
$\texttt{holonomy}(g_1) = \mathtt{true} = \texttt{holonomy}(g_2)$ and
$g_1 \ne g_2$.  Moreover, $g_2$ is the gauge transform of $g_1$ at vertex~0.
The Rashomon property is constructively witnessed with zero axioms.

\paragraph{Universal gauge invariance.}
The theorem \texttt{gauge\_preserves\_holonomy} is proved for \emph{all}
configurations, not just the two witnesses: for every $g \in \texttt{EdgeConfig}$,
$\texttt{holonomy}(\texttt{gaugeAt0}(g)) = \texttt{holonomy}(g)$.  The proof
proceeds by case analysis on all $2^3 = 8$ configurations.

\paragraph{The impossibility.}
By Theorem~\ref{thm:universal}, no description of a gauge field can
simultaneously be faithful (report the actual edge configuration), stable
(assign the same description to gauge-equivalent configurations), and decisive
(distinguish every pair of distinct configurations).  In the language of
physics: one cannot simultaneously specify a particular gauge potential,
demand gauge invariance of the specification, and retain the ability to
distinguish gauge-equivalent potentials.

\paragraph{Empirical illustration.}
On $\mathbb{Z}_2$ lattice gauge theory ($N \times N$ grids, $N = 4$--$10$),
gauge-variant link variables have within-orbit variance 0.25 (the Bernoulli
maximum) while gauge-invariant plaquette values have within-orbit variance
exactly~0.  The Rashomon set (gauge orbit) covers all possible link
configurations for a given holonomy.

Monte Carlo sampling on a $16 \times 16$ periodic lattice at 10 coupling
values ($\beta \in \{0.1, 0.2, \ldots, 2.0\}$; 2000 configurations per
$\beta$) confirms a coupling-dependent dose-response: the gauge coupling
$\beta$ controls the Rashomon set size.  Plaquette variance follows
$\Var(P) = \operatorname{sech}^2(\beta) / N^2$, decreasing monotonically
from $3.85 \times 10^{-3}$ at $\beta = 0.1$ to $2.71 \times 10^{-4}$ at
$\beta = 2.0$.  Wilson loop (2$\times$2) variance follows
$\Var(W) = 1 - \tanh^8(\beta)$, decreasing from ${\approx}1.00$ at weak
coupling to $0.23$ at strong coupling.  At weak coupling ($\beta \to 0$),
plaquettes are nearly uniform on $\{+1, -1\}$ and the configuration space is
maximally uncertain (large Rashomon set); at strong coupling
($\beta \to \infty$), plaquettes freeze to $+1$ and the configuration
concentrates near the ordered ground state (small Rashomon set).  Since 2D $\mathbb{Z}_2$ gauge theory is exactly solvable, the MC
serves as numerical verification of the analytic predictions rather than
an empirical discovery.  The contribution is the demonstration that the
gauge coupling continuously tunes the Rashomon set size---an exact
analogue of the collinearity parameter in the ML stability transition
experiment.

\paragraph{Resolution: gauge-invariant observables.}
The resolution in physics is well known: restrict attention to
\emph{gauge-invariant observables} such as the holonomy (Wilson loops),
field strengths, or scattering amplitudes.  In our framework, the holonomy
is the $G$-invariant projection: the gauge group acts on configurations,
and the holonomy is the unique function constant on each gauge orbit.
The resolution sacrifices decisiveness (one cannot recover the gauge
potential from the holonomy alone) to achieve stability and faithfulness.

\paragraph{Lean formalization.}
The Lean~4 types are \texttt{EdgeConfig} (a structure with fields
\texttt{e01 e12 e02 : Bool}), \texttt{holonomy} (the observation map),
and \texttt{gaugeSystem : ExplanationSystem EdgeConfig EdgeConfig Bool}.
The impossibility is \texttt{gauge\_impossibility} in
\texttt{GaugeTheory.lean}, derived with zero axioms.
